\uuid{bXTi}
\titre{Dérivée}
\auteur{bodin ; hanani ; mzari}
\organisation{exo7}
\licence{CC BY-NC-SA 4.0}
\datecreate{2020-01-20}
\niveau{L1}
\module{Algèbre}
\chapitre{Polynôme, fraction rationnelle}
\sousChapitre{Définition, degré, produit}
\difficulte{3}
\theme{Polynômes, dérivée}
\exo7qcmid{lille-1.polynomes.007}

\contenu{
\begin{question}
\qclassification{105.05}
\qtags{motcle=dérivée}
Soit $P(X) = \sum_{k=0}^n a_k X^k$. On associe le polynôme dérivé :
$P'(X) = \sum_{k=1}^n ka_k X^{k-1}$. Quelles sont les assertions vraies ?
\begin{answers}
    \good{Si $P$ est de degré $n\ge1$ alors $P'$ est de degré $n-1$.}

    \bad{Si $P'(X) = nX^{n-1}$ alors $P(X) = X^n$.}

    \good{Si $P'=P$ alors $P=0$.}

    \bad{Si $P'-Q'=0$ alors $P-Q=0$.}
\end{answers}
\begin{explanations}
C'est comme pour les primitives, il ne faut pas oublier la constante :
Si $P'=Q'$ alors $P=Q +c$.
\end{explanations}
\end{question}
}
